\section{Results}\label{sec:results}

We report BERT-Base in detail; DistilBERT behaves the same way unless stated.  Accuracies are means over chips, and hybrid placements keep 20\% of the
components in SLC unless stated otherwise.

\subsection{Cost of SVD and the Error-Free Chip}
\Cref{tab:svd} gives the accuracy before any hardware error. Post-training 8-bit quantisation costs
0.6--0.7 points. The SVD factorisation costs about two points even after fine-tuning; for DistilBERT,
fine-tuning slightly lowered accuracy, which suggests that 8 epochs over-fit. This cost is paid once and
is the same for every placement.

On a chip with no errors, every placement reaches 89.1--90.0\% on BERT-Base and 88.6--89.8\% on
DistilBERT, within about one point of the digital SVD model. The remaining differences come from
quantisation alone: each tier's crossbars have their own 8-bit scales, and the intermediate vector of
each layer is re-quantised to 8 bits, so the rounding depends on how the components are split between the
tiers.

\begin{table}[t]
\centering\footnotesize
\caption{SST-2 accuracy (\%) without hardware errors.}\label{tab:svd}
\setlength{\tabcolsep}{4pt}
\begin{tabular}{@{}lcccc@{}}
\toprule
Model & FP32 & INT8 & SVD & SVD, fine-tuned \\
\midrule
BERT-Base  & 92.43 & 91.86 & 89.33 & 90.14 \\
DistilBERT & 91.06 & 90.37 & 89.68 & 89.11 \\
\bottomrule
\end{tabular}
\end{table}

\subsection{Programming Noise}
\Cref{fig:noise} shows accuracy against programming noise. All-MLC storage fails first, falling to 66.3\% at
$\sigma=5\%$ and close to chance at 7.5\%. Keeping 20\% of the components in SLC by magnitude or Fisher stays
within about a point of all-SLC up to $\sigma=5\%$. The gradient placement keeps up until 5\% but falls behind
at higher noise (79.3\% against 83.5\% for magnitude at 7.5\%), and a random 20\% recovers little:
\emph{which} components are protected matters more than how many.

\begin{figure}[t]
\centering
\begin{tikzpicture}
\begin{axis}[accplot,xlabel={Programming noise $\sigma$ (\% of $G_{\max}$)},xmin=-0.3,xmax=10.3,
  xtick={0,1,2,3,5,7.5,10},legend pos=south west]
\addplot[pslc]  table[x=x,y=all_slc]{data/noise_bert.dat};   \addlegendentry{All SLC}
\addplot[pmag]  table[x=x,y=magnitude]{data/noise_bert.dat}; \addlegendentry{Magnitude}
\addplot[pfis]  table[x=x,y=fisher]{data/noise_bert.dat};    \addlegendentry{Fisher}
\addplot[pgrad] table[x=x,y=gradient]{data/noise_bert.dat};  \addlegendentry{Gradient}
\addplot[prand] table[x=x,y=random]{data/noise_bert.dat};    \addlegendentry{Random}
\addplot[pmlc]  table[x=x,y=all_mlc]{data/noise_bert.dat};   \addlegendentry{All MLC}
\end{axis}
\end{tikzpicture}
\caption{BERT-Base accuracy against programming noise, mean of 5 chips. Hybrids keep 20\% of the
components in SLC.}
\label{fig:noise}
\end{figure}

\subsection{Choosing and Sizing the Protected Set}
\Cref{fig:share} and \cref{tab:share} vary the share of components in SLC, with 20 chips per point.
Magnitude and Fisher lie on top of each other at every share. The gradient is close to them at
$\sigma=3\%$, but at 7.5\% it trails by 12 points with 5\% SLC and by 2--3 points at larger shares. Random
selection approaches the others only as the share grows towards one half.

\begin{figure}[t]
\centering
\begin{tikzpicture}
\begin{axis}[accplot,xlabel={Share of components in SLC (\%)},xmin=-2,xmax=102,
  xtick={0,10,20,30,50,100},ymin=50,legend pos=south east]
\addplot[pmag]  table[x=x,y=magnitude]{data/share_bert_075.dat}; \addlegendentry{Magnitude}
\addplot[pfis]  table[x=x,y=fisher]{data/share_bert_075.dat};    \addlegendentry{Fisher}
\addplot[pgrad] table[x=x,y=gradient]{data/share_bert_075.dat};  \addlegendentry{Gradient}
\addplot[prand] table[x=x,y=random]{data/share_bert_075.dat};    \addlegendentry{Random}
\end{axis}
\end{tikzpicture}
\caption{BERT-Base accuracy against the share of components in SLC at $\sigma=7.5\%$, mean of 20 chips.
0\% is all-MLC and 100\% all-SLC.}
\label{fig:share}
\end{figure}

\begin{table}[t]
\centering\footnotesize
\caption{BERT-Base accuracy (\%) by share of components in SLC, mean of 20 chips. 0\% and 100\% are
all-MLC and all-SLC.}\label{tab:share}
\setlength{\tabcolsep}{3.2pt}
\begin{tabular}{@{}llccccccc@{}}
\toprule
$\sigma$ & Placement & 0\% & 5\% & 10\% & 20\% & 30\% & 50\% & 100\% \\
\midrule
\multirow{4}{*}{5\%}
 & Magnitude & \multirow{4}{*}{66.8} & 88.0 & 88.6 & 88.7 & 88.9 & 88.9 & \multirow{4}{*}{88.0} \\
 & Fisher    & & 87.9 & 88.6 & 88.8 & 88.9 & 88.9 & \\
 & Gradient  & & 84.2 & 88.2 & 88.4 & 88.4 & 88.1 & \\
 & Random    & & 72.5 & 73.6 & 76.2 & 77.8 & 83.4 & \\
\midrule
\multirow{4}{*}{7.5\%}
 & Magnitude & \multirow{4}{*}{53.0} & 79.4 & 81.8 & 84.1 & 85.5 & 85.6 & \multirow{4}{*}{85.7} \\
 & Fisher    & & 80.0 & 82.2 & 84.2 & 85.1 & 85.8 & \\
 & Gradient  & & 67.5 & 79.0 & 81.2 & 82.4 & 83.2 & \\
 & Random    & & 53.6 & 54.2 & 55.2 & 57.5 & 67.3 & \\
\bottomrule
\end{tabular}
\end{table}

\textbf{Paired comparison.} Because all placements see the same chips, \cref{tab:paired} compares them
chip by chip against magnitude. Over all 300 chips of the study (3 noise levels, 5 SLC shares, 20 chips),
Fisher differs from magnitude by $+0.4\pm0.7$ sentences on BERT-Base and $-0.6\pm0.7$ on DistilBERT, out
of 872: no measurable difference at any noise level. The gradient is behind by $17.3\pm2.0$ and
$7.8\pm1.1$ sentences, and the gap grows with noise, to $40.4\pm5.0$ sentences on BERT-Base at
$\sigma=7.5\%$. \Cref{sec:discussion} explains why.

\begin{table}[t]
\centering\footnotesize
\caption{Paired difference to magnitude, in correctly classified sentences out of 872 (mean $\pm$
standard error over 100 chips per noise level, 300 overall).}\label{tab:paired}
\setlength{\tabcolsep}{2pt}
\resizebox{\columnwidth}{!}{%
\begin{tabular}{@{}llrrrr@{}}
\toprule
Model & Placement & $\sigma=3\%$ & 5\% & 7.5\% & All \\
\midrule
\multirow{3}{*}{BERT-Base}
 & Fisher   & $-0.1\pm0.5$ & $-0.1\pm0.7$ & $+1.4\pm1.9$ & $+0.4\pm0.7$ \\
 & Gradient & $-1.2\pm0.7$ & $-10.4\pm1.6$ & $-40.4\pm5.0$ & $-17.3\pm2.0$ \\
 & Random   & $-19.4\pm1.4$ & $-104\pm6$ & $-224\pm6$ & $-116\pm6$ \\
\midrule
\multirow{3}{*}{DistilBERT}
 & Fisher   & $-0.4\pm0.7$ & $-1.3\pm0.7$ & $+0.0\pm1.7$ & $-0.6\pm0.7$ \\
 & Gradient & $-0.9\pm0.7$ & $-4.0\pm1.0$ & $-18.4\pm2.6$ & $-7.8\pm1.1$ \\
 & Random   & $-11.2\pm1.1$ & $-66\pm5$ & $-186\pm7$ & $-88\pm5$ \\
\bottomrule
\end{tabular}}
\end{table}

\textbf{How much SLC is needed.} With magnitude ranking, the smallest share that comes within half a
point of all-SLC is 5\% at $\sigma=3\%$ for both models, 5\% (BERT-Base) to 10\% (DistilBERT) at
$\sigma=5\%$, and 30\% (BERT-Base) to 50\% (DistilBERT) at $\sigma=7.5\%$. In cells per weight
(\cref{sec:placement}), 10\% SLC needs 61\% and 30\% SLC 70\% of the cells of an all-SLC chip. The share
should therefore follow the expected noise: a small share suffices at moderate noise, and the saving
shrinks as noise grows.

\subsection{Stuck Cells}
\Cref{fig:fault} makes a share $p$ of all cells stuck, half at $G_{\max}$ and half at $G_{\min}$, with no
programming noise. Stuck cells are far more harmful than noise, because a stuck cell in a high-order slice
shifts its weight by a large fixed amount. All-MLC fails when one cell in a thousand is stuck. Up to
$p=0.2\%$, magnitude and Fisher with 20\% SLC match or exceed all-SLC, but from 0.5\% every placement
degrades quickly, and at 1\% even all-SLC loses 14--27 points across the two models. Survival then depends
on \emph{where} the stuck cells fall: at 0.5\% the worst of 20 BERT-Base chips reaches only 66.2\% with
magnitude placement. 

\begin{figure}[t]
\centering
\begin{tikzpicture}
\begin{axis}[accplot,xlabel={Stuck cells $p$ (\% of cells)},symbolic x coords={0,0.1,0.2,0.5,1,2,5},
  xtick=data,legend pos=north east]
\addplot[pslc]  table[x=x,y=all_slc]{data/fault_bert.dat};   \addlegendentry{All SLC}
\addplot[pmag]  table[x=x,y=magnitude]{data/fault_bert.dat}; \addlegendentry{Magnitude}
\addplot[pfis]  table[x=x,y=fisher]{data/fault_bert.dat};    \addlegendentry{Fisher}
\addplot[pgrad] table[x=x,y=gradient]{data/fault_bert.dat};  \addlegendentry{Gradient}
\addplot[prand] table[x=x,y=random]{data/fault_bert.dat};    \addlegendentry{Random}
\addplot[pmlc]  table[x=x,y=all_mlc]{data/fault_bert.dat};   \addlegendentry{All MLC}
\end{axis}
\end{tikzpicture}
\caption{BERT-Base accuracy against the share of stuck cells, no programming noise, mean of 20 chips.
The horizontal axis is not to scale.}
\label{fig:fault}
\end{figure}

\subsection{Other Non-Idealities}
\Cref{tab:chip} adds the remaining error sources one at a time. These experiments use evaluation batches of
64, so their error-free row differs slightly from \cref{fig:noise}.

\begin{table}[t]
\centering\footnotesize
\caption{BERT-Base accuracy (\%) under each of the other non-idealities. Hybrids keep 20\% of the components
in SLC.}\label{tab:chip}
\setlength{\tabcolsep}{2.4pt}
\begin{tabular}{@{}lccccc@{}}
\toprule
Condition & All SLC & All MLC & Magn. & Fisher & Grad. \\
\midrule
No errors & 89.7 & 90.4 & 90.6 & 89.6 & 90.3 \\
\midrule
Read noise $\sigma_r=2\%$ & 89.0 & 77.3 & 90.1 & 89.7 & 89.8 \\
Read noise $\sigma_r=5\%$ & 80.0 & 53.0 & 82.5 & 82.7 & 76.6 \\
$\sigma=5\%$, uniform levels & 88.2 & 67.3 & 88.9 & 89.0 & 88.3 \\
$\sigma=5\%$, middle levels $\times2$ & 88.2 & 62.5 & 88.4 & 88.4 & 87.8 \\
\midrule
1\% stuck, scattered & 61.3 & 50.7 & 56.4 & 55.6 & 52.6 \\
1\% stuck, $3\times3$ clusters & 66.7 & 50.8 & 54.0 & 53.1 & 51.5 \\
1\% stuck, bit lines & 57.0 & 50.4 & 50.4 & 52.2 & 51.0 \\
\midrule
ADC 7 bits & 90.0 & 89.2 & 89.3 & 89.7 & 90.1 \\
ADC 6 bits & 89.3 & 83.0 & 90.0 & 90.5 & 89.6 \\
ADC 5 bits & 51.1 & 51.1 & 51.1 & 51.6 & 53.0 \\
\midrule
IR drop $r=3\times10^{-4}$ & 88.9 & 89.8 & 90.4 & 91.1 & 90.5 \\
IR drop $r=10^{-3}$ & 87.0 & 90.6 & 88.1 & 88.3 & 88.6 \\
IR drop $r=3\times10^{-3}$ & 58.5 & 79.1 & 76.8 & 72.8 & 76.4 \\
\bottomrule
\end{tabular}
\end{table}

\textbf{Read and level-dependent noise} behave like programming noise. All-MLC fails (53.0\% at 5\% read
noise), magnitude and Fisher stay at or above all-SLC on BERT-Base , and the gradient trails them by six points at 5\% read noise.
Making the middle MLC levels twice as noisy costs all-MLC a further five points but the hybrids only about
half a point, since their most important components do not use MLC at all.

\textbf{Fault patterns.} At 1\% stuck cells, clusters harm all-SLC less than scattered faults, because
they damage fewer distinct weights, and broken bit lines harm it most; the hybrids are near chance for
every pattern.

\textbf{ADC resolution.} With the ADC range fixed by calibration, 6 bits suffice for all-SLC and for every
hybrid, while all-MLC needs 7. At 5 bits BERT-Base fails for every placement, and DistilBERT drops to
71--77\% (56.5\% for all-MLC). Below 6 bits results are not monotonic (BERT-Base: 51.1\% at 5 bits, 61.1\%
at 4), as most column sums round to a few levels; we count all of them as failures. With 5\% programming noise (DistilBERT, 3 chips), a 6-bit ADC keeps the magnitude hybrid at
88.5\%, against 56.0\% for all-MLC.

\textbf{ADC energy.} ADCs take 58\% of ISAAC's tile power~\cite{shafiee2016isaac}. Counting the
conversions of \cref{eq:adc} at 1.67\,pJ per 8-bit conversion~\cite{shafiee2016isaac,murmann2015adc},
with linear to exponential scaling in bits, the 6-bit BERT-Base hybrid needs 30\% fewer conversions than
all-SLC (23.0 against 33.0\,M per token) and 0.17--0.52$\times$ the ADC energy of an 8-bit all-SLC chip.

\textbf{IR drop.} At $r\leq3\times10^{-4}$ no placement loses more than a point. At $r=3\times10^{-3}$ IR
drop hurts SLC most (58.5\% against 79.1\% for all-MLC): by \cref{eq:irdrop} the highest conductances
lose the largest share of their current, and every stored SLC one sits at $G_{\max}$.

\begin{table*}[!t]
\centering\footnotesize
\caption{Accuracy (\%) under stuck cells, mean of 10 chips (best per column in bold). The all-SLC chip
stores every component in SLC and is the high-cost reference; the hybrid chip keeps 20\% of the
components in SLC, chosen by magnitude, and uses 34\% fewer cells. The chip column is the chip each
model is tested on; the training variants (\cref{sec:training}) differ only in the chip simulated during
training.}\label{tab:mapping}
\setlength{\tabcolsep}{5pt}
\renewcommand{\arraystretch}{1.1}
\begin{tabular}{@{}ll cccc c ccc c cc@{}}
\toprule
&& \multicolumn{4}{c}{Stuck cells, no noise} && \multicolumn{3}{c}{Stuck cells, $\sigma=5\%$} && \multicolumn{2}{c}{1\% stuck cells} \\
\cmidrule{3-6}\cmidrule{8-10}\cmidrule{12-13}
Chip & Mitigation & 0.2\% & 0.5\% & 1\% & 2\% && 0.2\% & 0.5\% & 1\% && clusters & lines \\
\midrule
\multicolumn{13}{@{}l}{\textbf{BERT-Base}} \\
\multirow{2}{*}{All SLC} & None & 85.4 & 80.1 & 60.9 & 51.1 && 85.6 & 78.4 & 64.6 && 66.5 & 57.2 \\
 & Fault-aware mapping & 89.6 & 83.7 & 75.8 & \textbf{55.9} && 86.1 & 82.4 & 76.5 && 81.3 & 77.5 \\
\arrayrulecolor{black!30}\specialrule{0.4pt}{1pt}{1pt}\arrayrulecolor{black}
\multirow{5}{*}{\shortstack[l]{Hybrid,\\20\% SLC}} & None & 88.5 & 76.6 & 55.9 & 50.1 && 84.9 & 70.3 & 52.1 && 53.7 & 50.1 \\
 & Fault-aware mapping & \textbf{90.1} & 85.6 & 65.1 & 52.2 && 87.4 & 80.4 & 58.1 && 64.3 & 56.9 \\
 & Noise-only training + fault-aware mapping & 89.8 & 87.0 & 67.6 & 51.7 && 88.0 & 82.7 & 58.0 && 65.8 & 58.9 \\
 & Noise + fault training + fault-aware mapping & 89.9 & \textbf{88.7} & 78.6 & 50.8 && \textbf{89.7} & 87.4 & 72.9 && 78.8 & 72.6 \\
\rowcolor{cmag!10}
 & Tier-aware training + fault-aware mapping & 88.8 & 88.3 & \textbf{82.9} & 53.0 && 89.1 & \textbf{87.8} & \textbf{80.5} && \textbf{85.0} & \textbf{77.7} \\
\midrule
\multicolumn{13}{@{}l}{\textbf{DistilBERT}} \\
\multirow{2}{*}{All SLC} & None & 88.1 & 85.4 & 73.7 & 57.9 && 86.9 & 81.7 & 68.0 && 79.6 & 71.1 \\
 & Fault-aware mapping & 88.0 & 87.6 & \textbf{84.8} & \textbf{62.3} && 87.6 & 86.8 & \textbf{80.7} && \textbf{85.6} & 81.5 \\
\arrayrulecolor{black!30}\specialrule{0.4pt}{1pt}{1pt}\arrayrulecolor{black}
\multirow{5}{*}{\shortstack[l]{Hybrid,\\20\% SLC}} & None & 87.1 & 77.9 & 62.4 & 50.2 && 85.4 & 74.9 & 59.5 && 55.9 & 54.0 \\
 & Fault-aware mapping & 88.4 & 85.1 & 67.0 & 53.3 && 87.7 & 83.0 & 65.1 && 68.5 & 68.0 \\
 & Noise-only training + fault-aware mapping & 87.3 & 84.5 & 68.2 & 51.5 && 85.9 & 82.8 & 59.9 && 66.1 & 65.9 \\
 & Noise + fault training + fault-aware mapping & \textbf{88.6} & \textbf{88.2} & 80.7 & 55.0 && \textbf{88.8} & 86.4 & 77.5 && 80.7 & 81.6 \\
\rowcolor{cmag!10}
 & Tier-aware training + fault-aware mapping & 88.2 & 88.1 & 84.7 & 57.4 && 88.2 & \textbf{86.9} & 79.8 && 83.4 & \textbf{84.4} \\
\bottomrule
\end{tabular}
\end{table*}

\subsection{Fault-Aware Mapping}
\Cref{tab:mapping} follows the stuck-cell results step by step, on 10 chips per point. 
We compare the hybrid against the all-SLC chip \emph{with} mapping, since mapping costs no hardware and
would be applied to any chip. Every row
sees the same chips, since the trained models keep the layer shapes. Differences quoted below are paired over these
chips, as mean $\pm$ standard error in accuracy points.

\textbf{Mapping helps in every setting.} Fault-aware mapping raises the hybrid at every fault rate and
pattern, for both models. At 0.5\% stuck cells BERT-Base rises from 76.6\% to 85.6\% and DistilBERT from
77.9\% to 85.1\%, and at 0.2\% the loss almost disappears. Chips also become more alike: at 0.5\% the
standard deviation over BERT-Base chips falls from 6.7 to 1.2 points. The largest gains come with
clustered and line faults (up to $13.9\pm3.3$ points), where moving important components away from a
damaged region matters most.

\textbf{Choosing the tier per chip does not help.} Joint tier and slot mapping (\cref{sec:mapping}) wins in
some settings and loses in others, by $-5.3\pm2.4$ to $+7.0\pm2.3$ points against fault-aware mapping, and
most differences are within two standard errors of zero. We attribute this to the cost model: its
second-order estimate underrates the large error of a stuck high-order cell, so it sometimes moves an
important component into MLC for a small modelled saving. The benefit of mapping comes from the choice of
slot, not of tier.

\textbf{Mapping alone is not enough.} The mapped hybrid matches all-SLC with mapping up to 0.2\% stuck
cells, but falls behind from 0.5\% upwards: at 1\% it reaches 65.1\% against 75.8\% on BERT-Base. In an
all-SLC chip every component sits in the tier where a stuck cell does least damage, and the mapping has
five times as many SLC slots to choose from.

\subsection{Noise- and Fault-Aware Training}
The last three rows of \cref{tab:mapping} add training to fault-aware mapping, and \cref{tab:trainnoise}
tests the trained models under programming noise alone.

\begin{table}[t]
\centering\footnotesize
\caption{Accuracy (\%) under programming noise alone, without stuck cells, mean of 10 chips. All hybrids
keep 20\% of the components in SLC, chosen by magnitude; $\sigma=0$ is an error-free chip.}\label{tab:trainnoise}
\setlength{\tabcolsep}{4pt}
\begin{tabular}{@{}lccc@{\hskip 10pt}ccc@{}}
\toprule
& \multicolumn{3}{c}{BERT-Base} & \multicolumn{3}{c}{DistilBERT} \\
\cmidrule(lr){2-4}\cmidrule(l){5-7}
Training & $\sigma=0$ & 5\% & 7.5\% & 0 & 5\% & 7.5\% \\
\midrule
No training       & 89.7 & 88.7 & 84.2 & 89.1 & 88.2 & 85.6 \\
Noise-only training & 89.9 & 89.7 & 86.3 & 87.4 & 87.6 & 85.6 \\
Noise + fault training & \textbf{90.9} & \textbf{89.9} & \textbf{88.8} & \textbf{89.9} & \textbf{89.3} & 87.6 \\
Tier-aware training & 89.4 & 89.3 & 88.3 & 87.8 & 88.4 & \textbf{88.0} \\
\bottomrule
\end{tabular}
\end{table}

\textbf{Training must include stuck cells.} Training with programming noise alone gives no consistent gain
under stuck cells: at 1\% it reaches 67.6\% and 68.2\%, against 65.1\% and 67.0\% with mapping alone.
Adding stuck cells to the training chip raises accuracy at 1\% to 78.6\% and 80.7\%, a gain of
$13.5\pm3.1$ and $13.7\pm2.6$ points.

\textbf{Training each tier on its own cells helps most.} Tier-aware training raises accuracy at 1\% stuck
cells further, to 82.9\% on BERT-Base and 84.7\% on DistilBERT ($+4.3\pm2.2$ and $+4.0\pm0.6$ points over
tier-unaware training), and to 80.5\% on BERT-Base when programming noise is added. Its advantage grows
with the density of errors: at 0.5\% the two training variants are equal, and at 0.2\% tier-aware training
is slightly behind. This is the price of training on the harsher hybrid chip, which costs 0.3 and 1.3
points on an error-free chip (\cref{tab:trainnoise}).

\textbf{Training and mapping are complementary.} Training makes the weights tolerate errors, while mapping
keeps the most sensitive values away from known faults. Each helps without the other: tier-aware training
without mapping reaches 76.1\% and 81.3\% at 1\% stuck cells, above mapping alone, and adding mapping still
gains $6.8\pm2.2$ and $3.5\pm1.2$ points.

\textbf{The gap to all-SLC closes.} With tier-aware training and mapping, the hybrid beats all-SLC with
mapping on BERT-Base at 0.5\% and 1\% stuck cells (88.3\% against 83.7\%, and 82.9\% against 75.8\%) and
with programming noise. On DistilBERT it matches all-SLC within one point up to 1\%, is ahead for broken bit
lines and two points behind for clusters. The hybrid does this with 34\% fewer cells. At 2\% stuck cells
every configuration still fails.

\textbf{Training also helps against programming noise.} Without stuck cells (\cref{tab:trainnoise}),
noise + fault training raises the hybrid at $\sigma=7.5\%$ from 84.2\% to 88.8\% on BERT-Base, and
tier-aware training from 85.6\% to 88.0\% on DistilBERT. Noise-only training helps BERT-Base (86.3\%) but
not DistilBERT. On an error-free chip, noise + fault training even gains about a point.
